\documentclass[12pt,a4paper]{article}
\usepackage[utf8]{inputenc}
\usepackage[french]{babel}
\usepackage[T1]{fontenc}
\usepackage{amsmath}
\usepackage{amsfonts}
\usepackage{amssymb}
\usepackage{graphicx}
\usepackage{hyperref}
\usepackage{listings}
\usepackage{xcolor}
\usepackage{geometry}

\geometry{margin=2.5cm}

\definecolor{codegreen}{rgb}{0,0.6,0}
\definecolor{codegray}{rgb}{0.5,0.5,0.5}
\definecolor{codepurple}{rgb}{0.58,0,0.82}
\definecolor{backcolour}{rgb}{0.95,0.95,0.92}

\lstdefinestyle{mystyle}{
    backgroundcolor=\color{backcolour},   
    commentstyle=\color{codegreen},
    keywordstyle=\color{magenta},
    numberstyle=\tiny\color{codegray},
    stringstyle=\color{codepurple},
    basicstyle=\ttfamily\footnotesize,
    breakatwhitespace=false,         
    breaklines=true,                 
    captionpos=b,                    
    keepspaces=true,                 
    numbers=left,                    
    numbersep=5pt,                  
    showspaces=false,                
    showstringspaces=false,
    showtabs=false,                  
    tabsize=2
}

\lstset{style=mystyle}

\title{Rapport Technique : Plateforme IoT NexusGuard \\ \large Village Mini - ENIF 6.0}
\author{Ingénierie NexusGuard}
\date{\today}

\begin{document}

\maketitle

\tableofcontents
\newpage

\section{Introduction}
Le projet \textbf{NexusGuard Village Mini} est une plateforme IoT complète conçue pour le monitoring et la gestion intelligente des ressources selon le paradigme \textbf{WEF Nexus} (Water-Energy-Food). Développée pour une exécution sur Raspberry Pi, cette solution permet d'assurer une irrigation autonome optimisée tout en protégeant les équipements matériels.

L'objectif principal est de maintenir un équilibre entre :
\begin{itemize}
    \item \textbf{Food (Sécurité Alimentaire) :} Maintenir l'humidité du sol pour la croissance des plantes.
    \item \textbf{Water (Gestion de l'Eau) :} Surveiller le niveau du réservoir et prévenir la marche à sec de la pompe.
    \item \textbf{Energy (Énergie/Protection) :} Surveiller la température ambiante pour éviter la surchauffe de la pompe.
\end{itemize}

\section{Architecture Matérielle}
Le système repose sur un micro-ordinateur \textbf{Raspberry Pi} connecté à une série de capteurs et d'actionneurs.

\subsection{Capteurs}
\begin{enumerate}
    \item \textbf{DHT22 :} Mesure de la température ambiante (°C) et de l'humidité de l'air (\%). Connecté au GPIO 17.
    \item \textbf{Capteur d'humidité du sol :} Connecté via un convertisseur analogique-numérique \textbf{ADS1115} sur le bus I2C (Canal 0).
    \item \textbf{HC-SR04 (Ultrason) :} Mesure la distance entre le haut du réservoir et la surface de l'eau pour calculer le volume restant (\%). Pins TRIG=23, ECHO=24.
\end{enumerate}

\subsection{Actionneurs}
Le système contrôle une \textbf{pompe à eau} via un relais (Logique inversée) connecté au GPIO 18. L'activation de la pompe dépend soit de la logique Nexus autonome, soit d'une commande manuelle via l'interface web.

\section{Architecture Logicielle}
La stack logicielle est moderne et optimisée pour le temps réel.

\subsection{Backend (Python/Flask)}
Le serveur est basé sur \textbf{Flask} et utilise \textbf{Flask-SocketIO} pour la communication bidirectionnelle en temps réel.
\begin{itemize}
    \item \texttt{app.py} : Gère les routes REST et les événements WebSocket.
    \item \texttt{main.py} : Contient la logique de lecture des capteurs et l'algorithme de décision Nexus.
    \item \texttt{config.py} : Centralise les paramètres (GPIO, seuils, intervalles).
\end{itemize}

\subsection{Frontend (Dashboard)}
L'interface utilisateur est développée en HTML5/CSS3 avec :
\begin{itemize}
    \item \textbf{Bootstrap 5} : Pour un design responsive et moderne (Dark Mode).
    \item \textbf{Chart.js} : Visualisation de l'historique des données sur les 30 dernières minutes.
    \item \textbf{Socket.io} : Mise à jour instantanée des indicateurs sans rechargement de page.
\end{itemize}

\section{Logique de Contrôle Nexus (WEF)}
L'algorithme de décision (dans \texttt{main.py}) vérifie trois conditions critiques avant d'autoriser l'irrigation :

\begin{equation}
Decision = (Soil < Threshold_{dry}) \land (Water > Threshold_{min}) \land (Temp < Threshold_{max})
\end{equation}

\subsection{Analyse des Dimensions}
\begin{itemize}
    \item \textbf{FOOD :} Si l'humidité du sol descend en dessous de 40\%, une demande d'eau est générée.
    \item \textbf{WATER :} Si le niveau d'eau est inférieur à 20\%, l'irrigation est bloquée pour protéger la pompe (protection contre la cavitation).
    \item \textbf{ENERGY :} Si la température dépasse 45°C, le système bloque la pompe par sécurité thermique.
\end{itemize}

\section{Interface Utilisateur}
Le dashboard propose deux modes de fonctionnement :
\begin{itemize}
    \item \textbf{Mode Automatique :} La logique Nexus contrôle totalement la pompe.
    \item \textbf{Mode Manuel :} L'utilisateur peut forcer l'état de la pompe via un bouton interactif.
\end{itemize}
Les seuils de déclenchement sont modifiables en direct via un formulaire de configuration rapide intégré à l'interface.

\section{Conclusion}
La plateforme NexusGuard offre une solution robuste et évolutive pour la gestion intelligente de micro-fermes ou de jardins connectés. L'intégration de la logique Nexus garantit non seulement la santé des cultures mais aussi la longévité du matériel de pompage.

\end{document}
